\section{Introduction}
\label{sec:intro}

\IEEEPARstart{W}{aste} can be recycled or composted only if the different materials are kept apart. In India this is a legal duty. The Solid Waste Management Rules, 2016 require every waste generator to segregate waste at source into bio-degradable, non-bio-degradable and domestic hazardous streams (Rule~4), and they prescribe green, white and black bins for bio-degradable, recyclable and other waste (Rule~15(h))~\cite{moefcc2016solid}. A camera that names the material of each item could help at a smart bin, on a sorting line or at a collection point, and machine learning is already common in solid waste management~\cite{abdallah2020artificial,fang2023artificial}.

Most of the vision papers we read solve an easier problem. They classify a photograph of one object, often on a plain background, as in TrashNet~\cite{yang2016classification} and the classifiers built on it~\cite{kaya2023classification}. A real heap, conveyor or street holds many items, some touching or half hidden, and each one has to be located as well as named. That is object detection. Majchrowska et al.~\cite{majchrowska2022deep}, whose work was our starting point, find litter with one network and classify each crop with a second. Multi-object waste detection therefore exists, but we found fewer such studies than single-object classifiers, and some detection studies still train on single-object studio images.

Detectors are also sensitive to training hyperparameters such as the learning rate, momentum, weight decay and augmentation strength. The TrashNet authors blamed the 22\% test accuracy of their convolutional neural network (CNN) on badly chosen hyperparameters~\cite{yang2016classification}. Metaheuristics have been used to tune CNNs, for example particle swarm optimization (PSO)~\cite{kennedy1995particle} for a mammography classifier~\cite{aguerchi2024cnn} and a beluga whale optimizer for an InceptionV3 waste classifier~\cite{sayed2024intelligent}. Two things are missing from the papers we read. First, the tuned waste models are classifiers; we found no waste detector whose training hyperparameters were tuned by a swarm optimizer. Second, the tuned model is compared with other metaheuristics or with nothing. We found no comparison with random search~\cite{bergstra2012random} at the same number of evaluations, and test results are single numbers without confidence intervals. With that evidence one cannot tell whether a gain comes from the optimizer or simply from trying many settings.

The data are the other weak point. Reviews report that many waste-sorting studies use simplified settings and artificially collected images~\cite{lu2022computer} that look different from real waste~\cite{fang2023artificial}. Accuracy also depends heavily on the dataset: the same YOLOv5x, trained and tested on each dataset separately, scored 84.9 AP$_{50}$ on PlastOPol but 63.3 on TACO~\cite{cordova2022litter}. A tuning gain measured on studio images may therefore say little about street photographs. For India the closest resource is the Dense Waste Segmentation Dataset (DWSD)~\cite{ali2025dwsd}. It reports segmentation pixel accuracy only, and we found no Indian real-world waste test set with published detection baselines.

In this work, which we call WasteSense, we put the tuning question to a controlled test. We train YOLOv8n~\cite{jocher2023ultralytics}, a detector small enough for a smart bin or a phone, and compare three ways (arms) of choosing six of its training hyperparameters: the library defaults (Arm~A), the best of 32 random samples (Arm~B), and our own PSO implementation with 8 particles over 4 rounds (Arm~C). Both searches spend 32 short proxy trainings. Each chosen setting is then trained once for 100 epochs and tested once on a held-out split, and every difference between arms gets a paired-bootstrap 95\% confidence interval. We did this twice. \genone used studio-style images in six classes. \gentwo added real-world photographs, a seventh class (OTHER) and an Indian test set built from DWSD that was never used in training. We ask three research questions:
\begin{enumerate}
\item[RQ1:] With an equal budget of 32 proxy trainings, does PSO find YOLOv8n training hyperparameters that beat random search and the defaults on a held-out test set?
\item[RQ2:] Does a short proxy training (12 epochs on 40\% of the data) rank hyperparameter settings the same way longer training does?
\item[RQ3:] Does the answer to RQ1 change when the data move from studio images to real-world and Indian scenes, and how large is the drop from studio to real-world?
\end{enumerate}

The contributions of the paper are:
\begin{itemize}
\item A small multi-object waste detector, YOLOv8n, trained on studio and real-world scenes with six classes in \genone and seven in \gentwo (Sections~\ref{subsec:datasets} and~\ref{subsec:detector}).
\item PSO tuning of six training hyperparameters of that detector: the learning-rate schedule, momentum, weight decay and two augmentation settings (Sections~\ref{subsec:space} and~\ref{subsec:pso}). We found no earlier swarm-tuned waste detector in the papers we reviewed.
\item A comparison of PSO, random search and the defaults at the same budget. The proxy was checked before use against a pass threshold set in advance (Spearman $\rho = 0.976$, $n = 8$, \genone only), the test split was used once, and every difference has a confidence interval (Sections~\ref{subsec:design}, \ref{subsec:proxy} and~\ref{subsec:eval}).
\item Separate results for studio, real-world and Indian test sets, including detection results on DWSD, for which we found no published detection baselines. Our DWSD boxes are derived from masks that merge touching items, and we had to infer the class mapping, so these scores partly reflect a labeling mismatch (Section~\ref{subsec:res-domain}).
\item The results as they came out, including the null ones. PSO did not beat random search in any way that matters. Tuning gave no gain over the defaults in \genone and a small one in \gentwo ($+0.017$ and $+0.020$ \mapfifty on the mixed test set), which did not carry over to real-world or Indian photographs. The code, search logs, result files and trained models are public.
\end{itemize}

Section~\ref{sec:related} reviews related work and states four research gaps (G1 to G4). Section~\ref{sec:methods} describes the methods and Section~\ref{sec:results} gives the results. Section~\ref{sec:discussion} discusses them and lists the limitations, Section~\ref{sec:future} covers future work, and Section~\ref{sec:conclusion} concludes.
